\section{Related Work}
\paragraph{Sequential recommendation.}
Transformer-based sequential models such as
\base~\cite{kang2018sasrec} and BERT4Rec~\cite{sun2019bert4rec}
(which adopts bidirectional masked prediction) are strong backbones for
capturing user behavior patterns.
GRU4Rec~\cite{hidasi2015gru4rec} pioneered RNN-based session
modeling, while SASRec introduced self-attention for next-item
prediction with causal masking. Recent efficient variants
include LightSANs~\cite{fan2021lightsans} with low-rank
attention decomposition for reduced complexity.

\paragraph{Text-enhanced recommendation.}
Item text features (from bag-of-words, neural encoders, or
\llm embeddings) provide complementary signals beyond
collaborative filtering. Early work used TF-IDF and
Item2Vec~\cite{barkan2016item2vec} for item representation;
CDL~\cite{wang2015cdl} jointly modeled text via stacked
denoising autoencoders. Recent approaches leverage pretrained
language models: UniSRec~\cite{hou2022unisrec} uses
BERT-encoded item text as universal transferable
representations, while VQ-Rec~\cite{hou2023vqrec} bridges
language and collaborative semantics via vector quantization.
Large language models have further expanded this line:
P5~\cite{geng2022p5} unifies recommendation in a text-to-text
paradigm, and TALLRec~\cite{bao2023tallrec} aligns LLMs with
recommendation via instruction tuning. However, these methods
typically encode item text with a single prompt or
representation; our work explores \emph{multi-view} \llm
embeddings via diverse prompt perspectives
(identity/function/audience/category) with explicit fusion and
alignment, and we show when this added complexity is justified
versus simpler alternatives.

\paragraph{Feature interaction and cross networks.}
Deep \& Cross Network (DCN)~\cite{wang2017dcn} and its
successor DCN-V2~\cite{wang2021dcnv2} explicitly model feature
crosses, complementing implicit interactions learned by DNNs.
Squeeze-and-Excitation Networks (SENet)~\cite{hu2018senet}
introduced channel-wise attention for adaptive feature
recalibration; FiBiNET~\cite{huang2019fibinet} adapted SENet
for CTR prediction with bilinear feature interactions. In
recommendation, cross networks have shown effectiveness in CTR
prediction; we adapt them for ID-text fusion in sequential
settings, and use SE-style blocks for per-view channel gating of \llm
features.

\paragraph{Multi-view learning in recommendation.}
Multi-view learning leverages complementary information from
different feature sources or modalities. Inspired by
CLIP~\cite{radford2021clip}, which aligns visual and textual
representations via contrastive learning, we design multi-view
prompts to elicit diverse semantics from language models. In
recommendation, multi-view approaches have been explored for
session-based settings~\cite{wang2023kgmvgnn} using
knowledge-enhanced graph neural networks. Our work applies
multi-view learning to \llm text embeddings, using different
prompt perspectives (identity, function, audience, category)
and whitening to decorrelate view representations before
fusion.

\paragraph{Contrastive learning in recommendation.}
Contrastive objectives such as \InfoNCE~\cite{oord2018cpc}
have been widely adopted for self-supervised representation
learning. In recommendation, contrastive methods align
user/item representations across views~\cite{wu2021cl4rec}.
We extend this to multi-view text-to-ID alignment with
cold-start reweighting.

\paragraph{Evaluation protocols.}
Sampled evaluation can mislead model
selection~\cite{krichene2020sampledmetrics}; full-ranking
evaluation provides more reliable comparisons but is
computationally expensive. We adopt full ranking with
stratified analysis to understand model behavior across
popularity strata.
